\section{Introduction}

Coding multi-agent systems (MAS) are a natural fit for software engineering tasks~\cite{swebench,multiagent2025,llmagents2024}.
A planner can inspect the issue and decide which files matter; an implementer can edit the code; a debugger can reason over failures; a reviewer can verify the final patch.
However, this decomposition creates a serving-system mismatch.
Each agent typically receives a different instruction prefix and different intermediate context, yet many agents repeatedly read the same repository files.
The shared content is therefore a large code-base segment embedded at different prompt positions, rather than a simple common prefix.

Modern LLM serving systems reduce prefill and decode cost through paged KV memory, structured program execution, semantic variables, workflow-aware cache scheduling, and KV-centric agent serving~\cite{kwon2023vllm,zheng2023sglang,parrot2024,kvflow2025,tokencake2025}.
However, existing prefix-cache mechanisms accelerate exact shared prefixes.
They do not directly address repeated middle or suffix spans.
A serving engine may observe that two prompts both contain the same large code file, but the cached KV tensors for that file were produced at different RoPE positions.
Naively reusing them can perturb attention behavior; using syntax or anchor similarity as a gate can be unsafe because the same AST shape or function name can hide behavior-changing edits.

This paper argues that coding MAS creates a structured opportunity for safe lossy KV reuse.
The word ``lossy'' in this work has a narrow meaning: KV tensors are reused across different prompt positions after RoPE delta alignment.
It does not mean that different or approximate code content may be reused.
Our safety rule is strict: AST, anchors, file paths, and span overlap are locators only; reuse is allowed only when the normalized code content signature is exactly identical.

We present \system, a template-guided serving stack for coding MAS.
\system connects three layers.
At the template layer, a workflow describes how Planner, Implementer, Debugger, and Reviewer agents will consume code-base segments.
At the scheduling layer, these templates emit coding-aware hints, including \prefetchhint, so a KVFlow-style engine can prepare code-base KV before a later agent needs it.
At the reuse layer, the prototype adapts cross-context KV reuse ideas from \kvcomm to code segments: it locates candidate spans, enforces \exactgate, and rotates cached K/V tensors using the RoPE position delta.
This framing borrows the original AgentTemplateKV insight that agent DAGs, roles, and tool edges are serving-time signals, but specializes it to repo-level coding workloads where large exact code files dominate the reusable context.

The contributions are:
\begin{itemize}[leftmargin=*]
\item We identify code-base repetition as a distinct cache-reuse pattern in coding MAS: large exact code segments recur across agents but not necessarily at prompt prefixes.
\item We design a template-derived code-base hint schema that exposes future code-base use through \prefetchhint while preserving a strict separation between scheduling hints and reuse safety gates.
\item We implement a code-specific exact-content reuse policy on top of an SGLang/KVFlow-style prototype, combining candidate localization, content-signature gating, and RoPE delta alignment without claiming KVFlow or \kvcomm themselves as new contributions.
\item We evaluate \system on SWE-bench Verified code-base workloads, controlled ablations, repo-level reuse experiments, and local pass@1 tests, showing zero false accepts for the full gate, near-lossless logit behavior, and cached-token/latency gains with small lossless-vs-lossy accuracy delta.
\end{itemize}
